% ============================================================================
% Chapter 5 — Discussion
% ============================================================================
\section{Discussion}
\label{sec:discussion}

\subsection{Theoretical contributions}
\label{sec:contributions}

DLVT contributes a closed-form, individual-level counterpart to the
organizational capability-trap literature. Where
\citet{repenning2002capability} and \citet{rahmandad2016capability} derive
bistable tipping structures---two basins, a separatrix, hysteresis---DLVT
shows that when the trap mechanism is routed through a single leader's
energetic budget, the structure collapses to one globally asymptotically
stable attractor: the trap is the default destination
(Theorem~\ref{thm:gas}). The practical difference is stark. In a bistable
world, prevention means staying in the healthy basin and rescue means a
large one-time push. In DLVT's world there is no healthy basin to defend;
only \emph{parameter} changes relocate the attractor, and
Corollary~\ref{cor:scope} says which parameters can do so. Relatedly, the
theory extends---rather than contests---human-capital theory
\citep{becker1993human} by adding the energetic-deployment constraint: the
wedge between capital owned and capital usable is the model's central
object, formalized in the energy-gated impact function \eqref{eq:impact}.

\subsection{Managerial implications}
\label{sec:implications}

Within the declared scope conditions, the model's implications invert a
common instinct. Scope reduction---delegating units, flattening
portfolios, trimming spans---feels like the natural remedy for executive
overload, but within the power-law family it is absorbed: capital
re-expands until the complexity load is restored (P3), and equilibrium
vitality is unchanged. What moves the equilibrium is the recovery side:
protected restoration capacity ($R$), slower capital decay for the same
impact ($\mu$), or higher conversion of impact into capital ($\alpha$),
consistent with evidence that respite interventions produce real but
decaying gains unless structurally maintained \citep{davidson2010sabbatical}
and with the recovery literature's emphasis on ongoing rather than episodic
restoration \citep{sonnentag2022recovery}. The model also cautions against
reading high-status incumbents' visible durability as evidence of
sustainability: at the low-vitality equilibrium, status ($C^{*}$) is at its
lifetime maximum precisely while deployable energy is chronically below
threshold.

\subsection{Limitations}
\label{sec:limitations}

Six limitations bound the theory's claims; we state them as bluntly as the
results.

\paragraph{Deterministic dynamics.} The model has no stochastic forcing.
Real careers experience shocks (health events, market dislocations,
reorganizations) that could produce excursions and regime misclassification
in finite samples; the global-stability result guarantees return to the
attractor but says nothing about excursion costs.

\paragraph{Single leader, closed loop.} Complexity is modeled as a
deterministic image of the leader's own capital
(Assumption~\ref{ass:scope}). This is a declared scope condition, not an
empirical claim: peer competition, organizational design, and market
turbulence all generate load that is orthogonal to own-capital
accumulation. Where exogenous load dominates, the exogenous-complexity
variant (Proposition~\ref{prop:band}) is the right object, and the
endogenous feedback results do not apply.

\paragraph{Illustrative calibration.} Table~\ref{tab:parameters} has no
empirical anchor. The honest summary of the baseline is that it sits near
its own regime boundary: $\mu/\alpha=2.0$ against a flip value of $2.163$
($8\%$ margin), and $P(\text{low-vitality}\mid\text{stable equilibrium})
\approx0.49$ across a $\pm2\times$ log-uniform hypercube. Qualitative
structure (uniqueness, stability, scope absorption) is
calibration-independent; regime labels are not.

\paragraph{Stipulated threshold.} The strategic threshold
$\Vstrat=0.5\,\Vmax$ (Definition~\ref{def:threshold}) is a labeling
convention. The baseline equilibrium sits only $5.9\%$ below it, and any
threshold below $0.47\,\Vmax$ reclassifies the baseline as sustainable.
Classification claims should therefore be read as conditional on $\theta$,
and empirical work should estimate the threshold rather than inherit it.

\paragraph{Identifiability.} The eleven parameters collapse to
approximately six effective dimensionless groups after nondimensionalizing
$V$ by $\Vmax$, $C$ by $\beta^{-1/\eta}$, and time by $\mu^{-1}$; within
the equilibrium predictions there are pronounced ridges---most importantly
in $\mu/\alpha$ (which pins $V^{*}$ via \eqref{eq:vstar}) and $R/\delta$
(which pins the flux threshold)---so cross-sectional equilibrium data alone
cannot identify the full parameter vector. Identification requires
trajectories, or shocks (below).

\paragraph{Level of formal guarantee.} The global-stability proof is
planar machinery (Dulac, Poincar\'e--Bendixson) and does not survive added
state dimensions. Extensions---teams, stochastic load, a second resource
stock---will need different tools, and we make no claims beyond two
dimensions.

\subsection{An empirical validation roadmap}
\label{sec:roadmap}

The theory is built to be tested, and the propositions of
\S\ref{sec:propositions} order the program from cheap to decisive.
First, \emph{measurement and trajectory fitting}: subjective vitality
\citep{ryan1997subjective} and recovery experience
\citep{sonnentag2007recovery} instruments administered longitudinally,
with career capital proxied by scope, compensation, and network centrality,
fitted with continuous-time structural equation models to test the damped
spiral signature and single-attractor convergence (P1). Second,
\emph{exogenous shocks to break the $\mu/\alpha$ degeneracy}: promotions,
reorganizations, and mandated leaves act as natural experiments that
separate depreciation from accumulation, which cross-sectional equilibrium
data cannot (P2, \S\ref{sec:limitations}). Third, the \emph{killer test}: a
pre-registered $2{\times}2$ field experiment crossing scope redesign
(coupling reduction) with recovery support (protected restoration
capacity) over a 12--18~month horizon---long enough for the model's
re-expansion dynamics to complete. DLVT predicts a main effect of recovery
support on equilibrium vitality, no durable main effect of scope redesign,
and conservation of the scope--capital product in the redesign arms (P3,
P6). A durable vitality gain from scope redesign alone, with the product
not conserved, falsifies the theory's central mechanism rather than a
calibration choice.

\subsection{Conclusion}
\label{sec:conclusion}

Leadership research has documented, at length, that success is energetically
expensive \citep{hambrick2005executive,rudolph2020healthy}. DLVT's
contribution is to show that a minimal, fully analyzable model of that
observation---capital builds complexity, complexity drains vitality,
vitality gates impact---delivers strong and falsifiable structure: one
equilibrium, globally stable, low-vitality by default under convex load,
indifferent to scope redesign within its kernel family, and movable only
from the recovery side. The corrected mathematics in this manuscript is
deliberately less dramatic than earlier drafts: no finite-time collapse, no
critical coupling, no fold catastrophe. What remains is, we believe, more
useful: a small set of sharp claims, each with a stated boundary and a
stated way to kill it.
